% chapters/50_neurosymbolic.tex
% Part of: AI / LLM / Agents / Hardware Notes

\section{Neurosymbolic AI \& Program Synthesis}

\subsection{Neural-symbolic integration}

\begin{intuitionbox}
\textbf{Connection to Berkeley Week 8:} The course frames this as the fundamental interface between \emph{learned/statistical reasoning} (neural) and \emph{structured/logical reasoning} (symbolic). Robust AI agents need both: neural systems to handle perception, ambiguity, and learning from data; symbolic systems to enforce constraints, enable compositionality, and provide interpretability. Neurosymbolic AI seeks to combine their complementary strengths.
\end{intuitionbox}

\textbf{Motivation for Integration:}
Pure neural systems excel at pattern recognition and learning from data, but suffer from lack of explicit reasoning, compositionality, and interpretability—they treat the learned function as a black box. Pure symbolic systems are logically sound and transparent, but are brittle (require perfect knowledge engineering) and don't learn from data or handle perceptual ambiguity.

Two key failure modes motivate integration:
\begin{itemize}
    \item \textbf{Neural hallucination:} Neural systems generate outputs that violate logical constraints or domain knowledge. A vision model might identify an object as ``cat'' despite strong spatial or contextual evidence to the contrary.
    \item \textbf{Symbolic brittleness:} Symbolic systems fail to handle noisy, ambiguous, or high-dimensional perceptual input. A rule-based reasoner requires crisp, discrete facts and cannot gracefully degrade with partial information.
\end{itemize}

\textbf{Key Integration Approaches:}

\begin{enumerate}
    \item \textbf{Learning + Reasoning Pipeline} (neural perception → symbolic reasoner)

    Neural network produces a learned representation or embedding. A symbolic reasoner then operates on this representation to reach conclusions.
    \begin{itemize}
        \item \textbf{Neuro-Symbolic Concept Learner (NS-CL):} Learn object attributes (color, size, shape) from visual input via a neural network, then apply symbolic logical rules to reason about relationships. The neural part is trained end-to-end by backpropagating through the symbolic reasoning layer.
        \item Advantage: reasoning is interpretable and provably correct once attributes are extracted; learning is data-driven.
    \end{itemize}

    \item \textbf{Differentiable Logic} (make symbolic operations learnable)

    Embed symbolic operations in a differentiable framework so they can be trained with gradient descent.
    \begin{itemize}
        \item \textbf{Differentiable Theorem Proving:} Represent logical inference as a differentiable computation graph. A learned ``attention weight'' can guide which axioms and rules to apply at each step, optimizing the proof search.
        \item \textbf{Neural Theorem Provers:} Use a neural network to suggest the next proof step in an interactive theorem prover; a symbolic verifier (e.g., Lean) checks correctness. Only valid steps are accepted, ensuring soundness.
        \item Advantage: combines the scalability and generalization of neural search with the soundness guarantee of symbolic verification.
    \end{itemize}

    \item \textbf{LLMs as Symbolic Reasoners} (generate symbolic programs for execution)

    Use an LLM to generate a symbolic program (e.g., logical rules, pseudocode, formal specification) that is then executed by a formal interpreter, which provides feedback.
    \begin{itemize}
        \item \textbf{Program-Aided Language Models (PAL):} Query the LLM to write Python code that solves a math or reasoning problem. The code is executed, and the interpreter output (correct or error message) serves as ground truth. Intermediate results guide the LLM's reasoning.
        \item \textbf{Tree of Thought with Symbolic Verifier:} Explore a tree of candidate solutions (generated by the LLM), using a symbolic verifier or test suite to evaluate each branch. Prune invalid branches; expand promising ones.
        \item Advantage: LLMs are powerful at code and language generation; symbolic execution provides deterministic, verifiable results.
    \end{itemize}

    \item \textbf{Constraint Satisfaction} (symbolic filter on neural outputs)

    Run a neural model (e.g., LLM) unconstrained, then apply symbolic constraints as a post-processing filter or during decoding.
    \begin{itemize}
        \item \textbf{Constrained Decoding:} Use an SMT solver or type checker to restrict the LLM's next-token distribution to only those tokens that maintain syntactic validity or satisfy domain constraints.
        \item Example: in generating a SQL query, constrain the decoder to only emit tokens that keep the query syntactically valid (balanced parentheses, valid column names).
        \item Advantage: leverages the LLM's flexibility while guaranteeing structural correctness.
    \end{itemize}
\end{enumerate}

\textbf{System 1 vs. System 2 Framing:}
A useful cognitive analogy: \textbf{System 1} (fast, intuitive, pattern-matching) maps to neural processing; \textbf{System 2} (slow, deliberate, rule-based reasoning) maps to symbolic processing. Neurosymbolic AI aims to combine both—rapid intuitive judgments refined by careful deliberate reasoning. This mirrors human cognition: we use pattern recognition for quick assessment, then apply logical rules when a decision matters.

\subsection{Program induction}

\textbf{Goal and Definition:}
Program induction is the task of learning a \emph{program}—a structured, executable representation with formal semantics—from examples, rather than learning a flat neural function. The inferred program should generalize beyond the training examples.

\textbf{Program Synthesis:} The broader term for generating code that satisfies a specification. Specifications can be input/output examples, natural language descriptions, or formal logical specifications.

\textbf{Key Approaches:}

\begin{enumerate}
    \item \textbf{Search-Based Synthesis}

    Enumerate candidate programs from a domain-specific language (DSL) and check which ones satisfy the specification.
    \begin{itemize}
        \item \textbf{FlashFill (Microsoft Excel):} User provides a few input/output examples. The system searches a DSL of string operations (substring, regex, concatenation) to find a program that matches all examples and generalizes.
        \item \textbf{PROSE Framework:} General-purpose synthesis engine; can learn string transformations, data wrangling rules, and more. Combines search with ranking heuristics to prioritize likely programs.
        \item Advantage: guaranteed correct (checks against examples); disadvantage: search space is exponential; works best for small, well-defined DSLs.
    \end{itemize}

    \item \textbf{Inductive Logic Programming (ILP)}

    Learn logical rules (Horn clauses) from positive and negative examples. A rule is ``good'' if it covers all positives and few negatives.
    \begin{itemize}
        \item Example: given facts (``parent(tom, bob)'', ``parent(bob, ann)'') and positive examples (``ancestor(tom, ann)''), infer the rule ``ancestor(X, Y) :- parent(X, Y) or ancestor(X, Z), parent(Z, Y).''.
        \item Advantage: results are human-readable and interpretable.
    \end{itemize}

    \item \textbf{Neural Program Synthesis}

    Train a neural model (e.g., a transformer or RNN) to generate programs.
    \begin{itemize}
        \item \textbf{Guided Search:} The neural model ranks or prioritizes candidate programs. Promising programs (those passing intermediate tests) are expanded; others are pruned.
        \item \textbf{Direct Generation:} The neural model directly generates a program as text (similar to a language model). Training uses execution feedback: if the generated program passes tests, the sequence is reinforced; if it fails, backtrack or adjust.
        \item Advantage: scales to large program spaces; learns patterns from large corpora of code.
    \end{itemize}

    \item \textbf{DreamCoder}

    An alternating two-step algorithm:
    \begin{itemize}
        \item \textbf{Step 1 (Synthesis):} Given a library of primitive functions, synthesize programs that solve the current task set.
        \item \textbf{Step 2 (Learning):} Analyze the synthesized programs to identify recurring substructures. Add the most useful substructures as new library functions (abstraction). These become available for future synthesis.
    \end{itemize}
    Iteration drives \emph{compression}—the library becomes more expressive and specialized to the domain, making future synthesis faster and more elegant. Inspired by cognitive development: humans learn building-block skills that enable solving more complex tasks.

    \item \textbf{Execution-Guided Synthesis}

    Generate a partial program, execute it to compute intermediate results, and use those results to guide generation of the next step. Critical for long-horizon code generation.
    \begin{itemize}
        \item Example: to synthesize a data wrangling program, first run the initial transformation (e.g., ``load CSV''), observe the intermediate table, then generate the next step (e.g., ``filter rows where column > 5'') with full visibility of the data shape.
        \item Advantage: avoids dead-end paths (can detect errors early); exploits concrete intermediate feedback.
    \end{itemize}
\end{enumerate}

\textbf{Program Induction vs. Synthesis:}
\begin{itemize}
    \item \textbf{Induction:} Infer a general rule or principle from examples (data-centric).
    \item \textbf{Synthesis:} Generate specific code that satisfies a specification (spec-centric).
\end{itemize}
The terms are often used interchangeably, but induction emphasizes learning abstractions from data, while synthesis emphasizes satisfying a constraint.

\subsection{AlphaCode \& code generation}

\textbf{AlphaCode (DeepMind, 2022):}

AlphaCode was the first LLM to achieve competitive-programmer level performance on Codeforces, reaching the top 54th percentile—roughly equivalent to a strong intermediate-level competitive programmer.

\begin{notebox}
For context: Codeforces is a platform hosting algorithmic programming contests. Problems range from simple implementation to complex combinatorial optimization. Success requires strong algorithmic knowledge, ability to handle edge cases, and skill in debugging.
\end{notebox}

\textbf{Key Innovations:}

\begin{enumerate}
    \item \textbf{Massive Code-Specific Pretraining}

    AlphaCode was pretrained on a large, carefully curated dataset of code (GitHub + Codeforces problem solutions and contest history). This focused domain expertise is crucial: general-purpose language models underperform on competitive coding.

    \item \textbf{Problem Understanding}

    Fine-tune a transformer encoder on problem statements (natural language descriptions + constraints). The encoder produces a rich embedding of the problem. The decoder is conditioned on this embedding when generating candidate solutions.

    \item \textbf{Large-Scale Sampling}

    Generate a very large number of candidate solutions per problem (1M+ samples). The key insight: even if the success rate is low (~0.1\%), quantity compensates. With 1M samples, you can expect hundreds of correct solutions.

    \item \textbf{Filtering and Clustering}

    Run all generated programs against public test cases (unit tests provided with the problem). Correct solutions are retained; incorrect ones discarded. Group the correct solutions into clusters by runtime behavior or structure. Sample one solution per cluster for final submission (avoids redundancy; improves diversity of submissions).
\end{enumerate}

\textbf{AlphaCode 2 (2023):}

Combined with Gemini, achieved top 15th percentile—a substantial improvement. Likely driven by the stronger underlying language model and continued scaling of compute.

\textbf{Code as Structured Text:}

A key insight distinguishes code generation from natural language generation:
\begin{itemize}
    \item \textbf{Code has formal semantics.} It can be executed and tested. A program either passes or fails unit tests; there is an unambiguous ground truth.
    \item \textbf{Natural language is subjective.} Multiple answers can be correct; evaluation is fuzzy and context-dependent.
\end{itemize}
This enables \emph{automated, unambiguous evaluation}: we can run generated code and immediately know if it works. This feedback signal is invaluable for iterative refinement.

\textbf{Connection to Coding Agents:}

Modern coding agents (Claude Code, Devin, SWE-agent, etc.) extend AlphaCode's ideas:
\begin{itemize}
    \item Generate code with an LLM.
    \item Execute the code; capture output or error messages.
    \item Feed errors back to the LLM to refine the solution.
    \item Repeat until the code passes all tests or the agent exhausts its attempts.
\end{itemize}
This is a form of \emph{execution-guided program synthesis}. The agent's loop is: generate → execute → error feedback → refine. Each iteration narrows the search space toward a correct solution. Unlike AlphaCode's offline filtering, agents can refactor and debug in real time.

\subsection{Formal verification + LLMs}

\textbf{Formal Verification Primer:}

Formal verification is the task of \emph{mathematically proving} that a program, circuit, or system satisfies a specification. Unlike testing (which checks examples), verification is exhaustive—it proves correctness for all possible inputs.

Key tools:
\begin{itemize}
    \item \textbf{Model Checking:} Exhaustively explore the state space of a system (e.g., finite automaton, hardware design). Tools: Spin, NuSMV.
    \item \textbf{Theorem Proving:} Construct a formal proof using axioms and inference rules. Tools: Coq, Lean, Isabelle.
    \item \textbf{SMT Solvers:} Decide satisfiability of logical formulas (Boolean, arithmetic, etc.). Tools: Z3, CVC5.
\end{itemize}

\textbf{LLMs for Formal Proofs:}

LLMs can assist in interactive theorem proving (ITP) by suggesting proof steps. The key is that an automated verifier checks each step; only valid steps are accepted, ensuring soundness.

\begin{warnbox}
\textbf{Critical Safety Property:} LLMs are unreliable at proof generation and often hallucinate invalid steps. However, when an LLM is tightly coupled with a \emph{verifier}, the verifier rejects invalid steps. The system as a whole is sound: every accepted step is mathematically correct, even if the LLM sometimes generates garbage.
\end{warnbox}

\textbf{Key Systems:}

\begin{enumerate}
    \item \textbf{Lean + LLM}

    Lean is an interactive theorem prover. A proof is built by applying tactics (e.g., ``intro'', ``simp'', ``rw'') to simplify the goal. An LLM (via agents or decision trees) suggests which tactic to apply next. Lean's kernel verifies the tactic is legal; if not, the tactic is rejected and the LLM tries again.

    \textbf{Hyper Tree Proof Search (HTPS):} A search algorithm that explores the proof tree using an LLM to evaluate and expand promising nodes. The verifier ensures only valid proof steps are recorded.

    \textbf{AlphaProof (DeepMind, 2024):} Combined tree search and LLM. Achieved remarkable results on International Mathematical Olympiad (IMO) geometry and algebra problems, with fully formal proofs generated in Lean. Proofs are not just correct by machine check; they are human-readable and insights can be extracted.

    \item \textbf{MathLib (Lean 4)}

    A large, curated library of formalized mathematics (thousands of theorems and definitions in Lean). Serves as: (1) a training corpus for proof-generating models, and (2) a foundation for new proofs (reuse existing theorems).

    \item \textbf{LLMs as Specification Writers}

    Before verifying a system, you must specify its desired behavior formally (e.g., in temporal logic, Z notation, or TLA+). LLMs can help translate natural language requirements into formal specs:
    \begin{itemize}
        \item Engineer writes: ``The system must allow any user to add items to their cart and checkout without losing data.''
        \item LLM generates a TLA+ spec that formalizes this constraint.
        \item A model checker verifies the system against the spec.
    \end{itemize}
    Advantage: LLMs accelerate the costly specification-writing step.
\end{enumerate}

\textbf{Constrained Decoding for Correctness:}

During code or proof generation, use SMT solvers, type checkers, or grammar constraints to restrict the LLM's output to syntactically and semantically valid tokens:
\begin{itemize}
    \item \textbf{Type Checking:} Only emit tokens that maintain type safety (useful in Lean or other typed languages).
    \item \textbf{SMT Constraints:} If the proof step must satisfy a logical constraint (e.g., the variable must be positive), constrain the decoder to only suggest steps that satisfy the constraint.
\end{itemize}
This dramatically reduces the LLM's hallucination rate by ruling out invalid outputs before they are generated.

\textbf{Limitations and Open Questions:}

\begin{itemize}
    \item \textbf{Hallucination Still Happens:} Even with verification, LLMs frequently generate invalid proof steps. The verifier rejects them, but the search process must backtrack, slowing convergence.
    \item \textbf{Exponential Search Complexity:} Proof search is exponential in problem difficulty. For very hard problems (like hard IMO problems), even with LLM guidance, the search space is enormous. Success requires careful problem reduction and domain-specific heuristics.
    \item \textbf{Generalization:} Most systems are trained on specific domains (IMO geometry, software verification). Generalization to new domains is unclear.
    \item \textbf{Human Insight:} The most impressive results (AlphaProof on IMO) combine LLM guidance with symbolic search and external tools (e.g., computational geometry for angle-chasing). Pure neural or pure symbolic approaches fall short.
\end{itemize}

